% !PRINTED 480
\weberpage{480}

Wählen wir, wie oben, die Grösse $\varrho$ als rationale Function der $\alpha$ so, dass die $\Pi(m)$ durch die Permutationen der $\alpha$ sich ergebenden Grössen
\[
\varrho, \varrho', \varrho'' \ldots
\]
alle von einander verschieden sind, so ist
\[
G(t) = (t - \varrho)(t - \varrho')(t - \varrho'') \ldots
\]
eine Function in $\Omega$. Nun ist jede von den Grössen $\varrho, \varrho', \varrho'' \ldots$ eine primitive Grösse des Normalkörpers $N = \Omega(\alpha, \alpha_1 \ldots \alpha_{m-1})$, und jede von ihnen ist also die Wurzel einer \gesperrt{Galois'schen} Resolvente $\mu^{\text{ten}}$ Grades. Je $\mu$ von diesen Grössen sind die Wurzeln von einer solchen Gleichung. Es muss also $G(t)$, was keine gleichen Wurzeln hat, in lauter irreducible Factoren $\mu^{\text{ten}}$ Grades zerfallen, und daraus ergiebt sich noch, \gesperrt{dass $\mu$ ein Theiler von $\Pi(m)$ ist. Zugleich ist $\mu$ der Grad der Galois'schen Gruppe} von $F(x)$.

Hat der Grad der \gesperrt{Galois'schen} Resolvente einer Gleichung $m^{\text{ten}}$ Grades den grössten Werth $\Pi(m)$, so sagen wir, mit einem von \gesperrt{Kronecker} herrührenden Ausdruck, die Gleichung hat \gesperrt{keinen Affect}. Sie hat einen um so höheren Affect, je niedriger der Grad $\mu$ der \gesperrt{Galois'schen} Resolvente ist. Den Quotienten $\Pi(m) : \mu$, der immer eine ganze Zahl, höchstens gleich $\Pi(m)$ und mindestens gleich 1 sein muss, wollen wir den \gesperrt{Grad des Affectes} nennen, der also bei einer Gleichung ohne Affect den Werth 1 hat. Wenn der Affect den möglichst hohen Grad $\Pi(m)$ hat, dann sind die Wurzeln der Gleichung selbst im Rationalitätsbereich $\Omega$ enthalten, die Gleichung also gelöst.

Wenn durch Adjunction einer algebraischen Grösse zu $\Omega$ die \gesperrt{Galois'sche} Resolvente reducibel wird, so entsteht eine neue Resolvente niedrigeren Grades, und der Grad des Affects der Gleichung $F(x) = 0$ vergrössert sich. Man nähert sich also dadurch der Lösung der Gleichung.

Die \gesperrt{Galois'sche} Auffassung der Aufgabe, eine Gleichung $F(x) = 0$ zu lösen, besteht darin, dass durch auf einander folgende Adjunction von algebraischen Grössen möglichst einfacher Art die Gruppe allmählich verkleinert oder der Affect erhöht werden soll, bis er seinen höchsten Grad erreicht hat.

\begin{quote}
\gesperrt{Die allgemeine Gleichung $m^{\text{ten}}$ Grades hat in dem Körper $\Omega$, der aus den rationalen Functionen der Co\"efficienten $a_1, a_2 \ldots a_m$ besteht, keinen Affect.}
\end{quote}
